\subsection{Anchor Quality Stratification on the Davis Kinase Benchmark}
\label{sec:anchor_quartile}

To understand how anchor binding strength affects prediction quality, we stratify the Davis kinase benchmark~\citep{davis2011kinase} by anchor pKi.
As described in Section~\ref{sec:cross_dataset}, anchors are retrieved from the DTC training set by Tanimoto chemical similarity, with a minimum anchor pKi of 7 (confirmed binders).
For the quartile analysis, we use the oracle anchor variant (strongest known binder per drug within Davis) to isolate the effect of anchor strength from retrieval quality; the main cross-dataset results in Table~\ref{tab:davis_cross_dataset} use realistic Tanimoto-retrieved anchors throughout.

\paragraph{Data characteristics.}
The Davis benchmark is a dense kinase selectivity panel: 442 kinases screened against 68 compounds, with binding measured as dissociation constant K$_d$.
The resulting pKd distribution is highly skewed: 91.4\% of interactions have pKd $< 7$ (non-binders at the assay sensitivity floor of pKd $= 5.0$), with only 2,500 of 29,170 pairs qualifying as true binders (pKd $\geq 7$).
This extreme class imbalance makes Davis a stringent test of binding discrimination---models must distinguish the 8.6\% true binders from a large background of non-binding kinases.

The oracle anchor pKi distribution is narrow (mean $9.10 \pm 0.70$, range 7.0--10.8) because each drug's strongest binder is typically a high-affinity kinase.
Anchor reuse is heavy: only 85 unique anchor proteins serve 29,170 interaction pairs, with a single kinase (DDR1) acting as anchor for 1,766 pairs (6.1\%).
This reflects the kinase family's shared structural fold---a small number of promiscuous kinases bind most compounds strongly.

Figure~\ref{fig:davis_distribution} shows the target pKi distribution, anchor pKi distribution, binder fraction per quartile, and interaction counts per quartile.

\begin{figure}[t]
    \centering
    \includegraphics[width=\textwidth]{fig_davis_distribution.pdf}
    \caption{Davis benchmark data characteristics. \textbf{(A)}~Target pKi distribution: 91.4\% of interactions are non-binders (pKi $< 7$), concentrated at the assay floor of pKd $= 5.0$. \textbf{(B)}~Anchor pKi distribution (anchors $\geq 7$ only): most anchors have pKi 8--10. \textbf{(C)}~Binder fraction per anchor quartile: Q3 has the highest binder rate (13.2\%) despite not having the strongest anchors. \textbf{(D)}~Interaction counts per quartile, with the number of unique anchor proteins annotated.}
    \label{fig:davis_distribution}
\end{figure}

\paragraph{Quartile definition.}
We partition the 29,170 interactions into four quartiles by anchor pKi: Q1 (7.0--8.7, $n = 8{,}404$, 41 anchors), Q2 (8.7--9.1, $n = 6{,}189$, 23 anchors), Q3 (9.1--9.5, $n = 7{,}513$, 27 anchors), and Q4 (9.6--10.8, $n = 7{,}064$, 20 anchors).
The binder fraction increases from Q1 (4.1\%) to Q3 (13.2\%) before declining slightly in Q4 (11.3\%), reflecting the fact that drugs with very strong anchors (Q4) tend to be highly selective compounds that bind fewer kinases overall.

Table~\ref{tab:davis_quartile_dist} summarizes the target pKi statistics per quartile.

\begin{table}[t]
    \centering
    \caption{Davis target pKi distribution stratified by oracle anchor pKi quartile. Anchors restricted to pKi $\geq 7$. Binders defined as pKi $\geq 7$.}
    \label{tab:davis_quartile_dist}
    \begin{tabular}{lrrccccc}
        \toprule
        Quartile & $n$ & Binders & \% Bind & Median & Mean & Std \\
        \midrule
        Q1 [7.0--8.7] & 8,404  & 347   & 4.1  & 5.00 & 5.27 & 0.65 \\
        Q2 [8.7--9.1] & 6,189  & 359   & 5.8  & 5.00 & 5.34 & 0.76 \\
        Q3 [9.1--9.5] & 7,513  & 994   & 13.2 & 5.00 & 5.67 & 1.02 \\
        Q4 [9.6--10.8] & 7,064 & 800   & 11.3 & 5.00 & 5.59 & 1.07 \\
        \midrule
        Overall & 29,170 & 2,500 & 8.6 & 5.00 & 5.46 & 0.90 \\
        \bottomrule
    \end{tabular}
\end{table}

\paragraph{Performance comparison.}
We compare three models: (1)~\textbf{V2-650M}, anchor transfer with ESM-2 650M protein embeddings (1,280-dim); (2)~\textbf{V2-35M}, anchor transfer with ESM-2 35M embeddings (480-dim); and (3)~\textbf{DeepDTA}~\citep{ozturk2018deepdta}, a pairwise CNN model that encodes the target protein's amino acid sequence directly without any anchor mechanism.
All models are trained on Drug Target Commons (DTC) with protein-level 80/10/10 splits.

Table~\ref{tab:davis_quartile_perf} and Figure~\ref{fig:davis_quartile_performance} report CI, RMSE, and AUROC (binder threshold pKi $\geq 7$) per quartile.

\begin{table}[t]
    \centering
    \caption{Davis performance stratified by oracle anchor pKi quartile. CI~=~concordance index ($\uparrow$), RMSE in pKi units ($\downarrow$), AUROC with binder threshold pKi~$\geq$~7 ($\uparrow$). Best value per metric bolded. 95\% bootstrap CI in parentheses (1,000 resamples).}
    \label{tab:davis_quartile_perf}
    \begin{tabular}{lrccccccccc}
        \toprule
        & & \multicolumn{3}{c}{CI $\uparrow$} & \multicolumn{3}{c}{RMSE $\downarrow$} & \multicolumn{3}{c}{AUROC $\uparrow$} \\
        \cmidrule(lr){3-5} \cmidrule(lr){6-8} \cmidrule(lr){9-11}
        Quartile & $n$ & V2-650 & V2-35 & DTA & V2-650 & V2-35 & DTA & V2-650 & V2-35 & DTA \\
        \midrule
        Q1 & 8,404  & \textbf{0.691} & 0.583 & 0.517 & \textbf{0.786} & 1.102 & 1.077 & \textbf{0.750} & 0.614 & 0.487 \\
        Q2 & 6,189  & \textbf{0.706} & 0.585 & 0.450 & \textbf{0.896} & 0.940 & 1.155 & \textbf{0.766} & 0.693 & 0.479 \\
        Q3 & 7,513  & 0.649 & 0.643 & 0.510 & 1.047 & \textbf{0.958} & 1.292 & 0.726 & \textbf{0.706} & 0.513 \\
        Q4 & 7,064  & 0.706 & \textbf{0.744} & 0.514 & 0.940 & \textbf{0.984} & 1.271 & 0.806 & \textbf{0.832} & 0.501 \\
        \midrule
        Overall & 29,170 & \textbf{0.680} & 0.624 & 0.514 & \textbf{0.919} & 1.004 & 1.199 & \textbf{0.757} & 0.698 & 0.518 \\
        \bottomrule
    \end{tabular}
\end{table}

\begin{figure}[t]
    \centering
    \includegraphics[width=\textwidth]{fig_davis_quartile_performance.pdf}
    \caption{Davis performance by oracle anchor quartile. \textbf{(A)}~Concordance index: V2-650M dominates in Q1--Q2 (weak anchors), while V2-35M overtakes in Q4 (strong anchors). DeepDTA is near-random throughout. \textbf{(B)}~RMSE: V2-650M achieves the lowest error in Q1 (0.786) where the anchor signal is weakest. \textbf{(C)}~AUROC with binding threshold pKi~$\geq$~7: both V2 variants substantially outperform DeepDTA, which fails to exceed random (0.5) in any quartile.}
    \label{fig:davis_quartile_performance}
\end{figure}

\paragraph{Per-protein metric distributions.}
To examine whether the global improvements reflect consistent per-protein gains or are driven by a few outlier proteins, Figure~\ref{fig:davis_per_protein} shows violin plots of per-protein CI, RMSE, and AUROC across all 442 Davis kinases.
Table~\ref{tab:davis_per_protein} reports summary statistics.

\begin{table}[t]
    \centering
    \caption{Per-protein metric summary across 442 Davis kinases (oracle anchors). Each protein is evaluated independently over its 66-drug panel. V2-650M achieves the highest mean CI and AUROC; V2-ProstT5 achieves the lowest mean RMSE.}
    \label{tab:davis_per_protein}
    \begin{tabular}{lcccccc}
        \toprule
        & \multicolumn{2}{c}{CI $\uparrow$} & \multicolumn{2}{c}{RMSE $\downarrow$} & \multicolumn{2}{c}{AUROC $\uparrow$} \\
        \cmidrule(lr){2-3} \cmidrule(lr){4-5} \cmidrule(lr){6-7}
        Model & Mean & Median & Mean & Median & Mean & Median \\
        \midrule
        V2-650M  & \textbf{0.645} & \textbf{0.651} & 0.888 & 0.853 & \textbf{0.691} & \textbf{0.706} \\
        V2-35M   & 0.604 & 0.594 & 0.985 & 0.963 & 0.663 & 0.665 \\
        V2-ProstT5 & 0.610 & 0.603 & \textbf{0.831} & \textbf{0.809} & 0.635 & 0.639 \\
        ConPlex  & 0.567 & 0.581 & 5.720 & 5.701 & 0.535 & 0.556 \\
        ESM-DTA  & 0.496 & 0.496 & 1.311 & 1.274 & 0.499 & 0.485 \\
        DeepDTA  & 0.507 & 0.513 & 1.174 & 1.160 & 0.487 & 0.496 \\
        \bottomrule
    \end{tabular}
\end{table}

\begin{figure}[t]
    \centering
    \includegraphics[width=\textwidth]{fig_davis_per_protein_distributions.pdf}
    \caption{Per-protein metric distributions on the Davis benchmark (442 kinases). Violin plots show the full distribution; red lines indicate means, black lines medians. \textbf{(A)}~CI: V2-650M has the highest median (0.651) and tightest distribution. DeepDTA and ESM-DTA cluster at random (0.5). \textbf{(B)}~RMSE: ConPlex has an extreme scale mismatch (mean 5.72); V2-ProstT5 achieves the lowest RMSE (0.831). \textbf{(C)}~AUROC (binder~$\geq$~7~pKi): V2-650M's distribution is shifted rightward; pairwise baselines center at random.}
    \label{fig:davis_per_protein}
\end{figure}

\paragraph{Interpretation.}
Three findings emerge from the anchor-stratified analysis:

\textit{1.~Pairwise models fail entirely on out-of-distribution kinases.}
DeepDTA achieves per-protein CI~$= 0.507$ and AUROC~$= 0.487$---indistinguishable from random.
ESM-DTA performs similarly (CI~$= 0.496$, AUROC~$= 0.499$).
Both models encode the target protein's sequence directly, and because the Davis kinases are absent from the DTC training set, the sequence encoder has no mechanism to transfer knowledge from related proteins.
ConPlex achieves slightly better CI (0.567) through its contrastive ESM-2 protein embedding, but its RMSE is catastrophic (5.72) because its scoring function is not calibrated to predict pKi values.
These results confirm that pairwise sequence-based models are fundamentally limited for cross-dataset generalization.

\textit{2.~The 650M encoder compensates for weak anchors.}
V2-650M's advantage over V2-35M is largest in Q1 and Q2 (weak anchors, pKi 7.0--9.1): CI 0.691 vs.\ 0.583 in Q1, and 0.706 vs.\ 0.585 in Q2.
When the anchor provides only a weak binding signal, the larger protein encoder extracts more discriminative information from the protein representation to compensate.
In Q1, V2-650M also achieves the lowest RMSE of any quartile (0.786), demonstrating accurate value prediction even with weak anchor reference points.

\textit{3.~Strong anchors equalize encoder capacity.}
In Q4 (strong anchors, pKi 9.6--10.8), V2-35M overtakes V2-650M (CI 0.744 vs.\ 0.706, AUROC 0.832 vs.\ 0.806).
When the anchor provides a high-affinity reference point, even a smaller encoder can exploit the relative binding signal effectively---the anchor's strong binding to the drug constrains the prediction space sufficiently that additional encoder capacity yields diminishing returns.
This suggests that anchor quality and encoder capacity are partially substitutable: a strong anchor compensates for a weaker encoder, and a strong encoder compensates for a weaker anchor.

\textit{4.~ProstT5 structure-aware embeddings improve calibration but not ranking.}
V2-ProstT5 achieves the lowest per-protein RMSE (0.831, lower than V2-650M's 0.888) but lower CI (0.610 vs.\ 0.645).
ProstT5's structure-aware embeddings produce better-calibrated absolute pKi predictions, while ESM-2 650M's sequence embeddings produce better relative rankings.
This trade-off suggests that combining both embedding types could improve both calibration and discrimination.

\subsection{Ablation: Cosine Similarity vs.\ Learned Representations}
\label{sec:cosine_ablation}

A natural question is whether anchor transfer simply exploits protein similarity---predicting that proteins similar to a known binder will also bind---or whether the model learns drug-specific binding patterns beyond raw similarity. To answer this, we compare the V2-650M model's predictions against a simple baseline: the cosine similarity between the anchor and query protein embeddings, used directly as a predictor of binding affinity.

Table~\ref{tab:cosine_ablation} reports per-protein AUROC and CI for cosine similarity predictors using three embedding spaces, compared against the full V2-650M model.

\begin{table}[t]
    \centering
    \caption{Cosine similarity between anchor and query protein embeddings as a standalone binding predictor on Davis, compared to the full V2-650M model. Per-protein metrics averaged over 442 kinases. Cosine similarity captures some binding signal but substantially underperforms the learned model, which additionally encodes drug structure and three-way interactions.}
    \label{tab:cosine_ablation}
    \begin{tabular}{lcccc}
        \toprule
        Predictor & AUROC $\uparrow$ & CI $\uparrow$ & Spearman $\rho$ & Binder--non-binder $\Delta\cos$ \\
        \midrule
        \textbf{V2-650M (full model)} & \textbf{0.691} & \textbf{0.645} & --- & --- \\
        \midrule
        Cosine (ProstT5) & 0.658 & 0.582 & 0.146 & 0.068 \\
        Cosine (ESM-2 650M) & 0.574 & 0.552 & 0.141 & 0.028 \\
        Cosine (ESM-2 35M) & 0.545 & 0.539 & 0.083 & 0.015 \\
        \bottomrule
    \end{tabular}
\end{table}

\paragraph{Similarity is necessary but not sufficient.}
Cosine similarity between anchor and query embeddings does carry binding-relevant signal: proteins that bind the same drug tend to have more similar embeddings than non-binders ($\Delta\cos = 0.028$ for ESM-650M, $0.068$ for ProstT5). This signal is statistically significant (Spearman $\rho = 0.14$, $p < 10^{-100}$) and produces above-random per-protein AUROC (0.574 for ESM-650M, 0.658 for ProstT5). However, cosine similarity alone is far weaker than the full model: V2-650M achieves AUROC 0.691 and CI 0.645, representing a 20\% relative improvement over the best cosine baseline (ProstT5, AUROC 0.658).

\paragraph{The model learns more than similarity.}
The V2-650M model's predictions correlate only weakly with the cosine similarity of its own embedding space (Spearman $\rho = 0.257$ between V2 predictions and ESM-650M cosine). This means 74\% of the model's predictive variance is orthogonal to raw protein similarity. The three pairwise interaction modules---anchor--drug, query--drug, and anchor--query---capture drug-specific binding patterns that pure protein similarity cannot encode. In a per-protein comparison, V2-650M outperforms ESM-650M cosine similarity on 361 of 442 proteins (82\%), ties on 20, and loses on only 61.

\paragraph{Structure-aware embeddings improve the similarity signal.}
ProstT5 embeddings, which encode 3Di structural alphabet features from Foldseek \cite{van2024fast}, produce the strongest cosine similarity baseline (AUROC 0.658, $\Delta\cos = 0.068$). This is consistent with ProstT5's superior performance in the anchor transfer framework (V2-ProstT5-BDB global AUROC 0.768): structure-aware embeddings better capture the structural similarity between anchor and query proteins that underlies binding knowledge transfer.

\paragraph{Implications.}
This ablation demonstrates that anchor transfer is not reducible to protein similarity matching. The model leverages similarity as one input signal but learns additional drug-specific and interaction-specific patterns through its three-way architecture. The anchor mechanism provides a structured comparison framework---``how does this query compare to a known binder for this specific drug?''---that goes beyond asking ``how similar are these two proteins?'' This distinction is important because it means the model can discriminate between drugs even when the anchor--query similarity is constant, by learning how different chemical scaffolds interact differently with the anchor--query protein pair.

\subsection{Ablation: Oracle Anchor Selection}
\label{sec:oracle_anchors}

The main cross-dataset results (Table~\ref{tab:davis_cross_dataset}) use realistic anchors retrieved by Tanimoto similarity after excluding canonical duplicates. As an ablation, we compare against oracle anchor selection, where the strongest known binder of each drug \emph{within the Davis evaluation set itself} serves as anchor---providing an upper bound that requires privileged access to evaluation data.

With oracle anchors (restricted to pKi~$\geq 7$, $n = 29{,}170$), V2-650M achieves per-protein CI 0.645 and AUROC 0.691. This is slightly higher than the realistic protocol (CI 0.642, AUROC 0.669), as expected since oracle anchors select the globally strongest binder rather than the nearest training-set analog. The gap is modest ($\Delta$CI $= 0.003$, $\Delta$AUROC $= 0.022$), confirming that Tanimoto retrieval produces anchors of comparable quality to oracle selection for this benchmark. Pairwise baselines remain near random under both conditions.
